\chapter{The study}
\label{ch:study}

This chapter sets out the study the system was built for: what would be asked
of the participants, what would be measured, and how the data would be
analysed. None of it has been run. Stating it in full is what makes
\cref{ch:participant} an empty container rather than a gap.

\section{The question}

Does letting the machine take over part of the work help, and who does it
help?

The second half is the part that is usually left out. In a two-person
arrangement, an improvement for the driver is not automatically an improvement
for the wearer, and the existing evidence suggests it may be the
opposite~\cite{zhou2026proxemics}. The study is therefore designed so that
both people are measured on every trial, and so that a result can come out
mixed.

\section{Design}

\begin{figure}[htbp]
  \centering
  % Four tasks by three conditions, within each pair of participants.
\begin{tikzpicture}[x=1mm, y=1mm, font=\scriptsize]

  \foreach \j/\cond in {0/{C1 direct\\ teleoperation},
                        1/{C2 assisted\\ teleoperation},
                        2/{C3 spoken or typed\\ goal}}
    { \pgfmathsetmacro{\y}{-\j*14}
      \node[sense, text width=26mm, minimum height=11mm] at (-2,\y) {\cond}; }

  \foreach \i/\task/\note in {
      0/{T1 reaching}/{six targets, no grasping},
      1/{T2 pick and place}/{four blocks to matching mats},
      2/{T3 two-arm carry}/{one object, both grippers},
      3/{T4 hand to wearer}/{the robot serves the person carrying it}}
    { \pgfmathsetmacro{\x}{28 + \i*27}
      \node[stage, text width=24mm, minimum height=11mm] at (\x,14) {\textbf{\task}\\ \note}; }

  \foreach \j in {0,1,2}{
    \foreach \i in {0,1,2,3}{
      \pgfmathsetmacro{\x}{28 + \i*27}
      \pgfmathsetmacro{\y}{-\j*14}
      \pgfmathtruncatemacro{\blocked}{ ( \i==1 && \j<2 ) || ( \i==2 ) ? 1 : 0 }
      \ifnum\blocked=1
        \node[absent, text width=24mm, minimum height=11mm] at (\x,\y) {not
              currently possible};
      \else
        \node[stage, text width=24mm, minimum height=11mm, fill=green!6] at (\x,\y)
             {3 trials\\ per dyad};
      \fi } }

  \node[tag, text width=118mm, align=left] at (55,-50)
    {Orange cells are not a design choice. Grasping under direct teleoperation
     is not achievable on this platform, and the two-arm carry needs a grip
     wider than the hand opens. Both are measured results reported in
     \cref{ch:results}, and both are stated here so that the empty cells in
     \cref{ch:participant} are not read as missing data.};

  \node[tag, text width=118mm, align=left] at (55,-62)
    {Sixteen pairs of participants. Every pair sees every available cell.
     Condition order is balanced across pairs so that no condition is
     systematically first or last, and so that each condition follows every
     other equally often.};

\end{tikzpicture}

  \caption[The study design]{Three ways of driving, crossed with four tasks,
  within each pair of participants. Cells marked as not currently possible are
  measured limits of the platform rather than omissions, and are explained in
  \cref{ch:results}.}
  \label{fig:design}
\end{figure}

Sixteen pairs of participants, each pair doing every available combination of
condition and task, three trials each. Everything varies within a pair rather
than between pairs, because the differences between two people are larger than
the differences the study is looking for.

Condition order is balanced using a scheme that balances not only how often
each condition appears in each position, but also which condition follows
which. Simple counterbalancing leaves the second kind of imbalance in place,
and with a task that people get better at, it matters.

Two roles, not one. Each participant does one role for the whole session.
Swapping them mid-session would halve the sample and confound role with
fatigue, and the wearer's role involves carrying seventeen kilograms, which
nobody should do twice in one visit.

\section{The tasks}

\emph{T1, reaching.} Six marked targets, three per arm, reached in a given
order. No grasping. This is the only task available in all three conditions
and is therefore the one that carries the main comparison. Difficulty is
indexed in the standard way from the size of the target and the distance to
it~\cite{fitts1954information}, so trials of different geometry can be pooled.

\emph{T2, pick and place.} Four coloured blocks on a table, two coloured mats,
each block to the mat of its own colour. This is the task the autonomous mode
performs end to end from a typed sentence.

\emph{T3, two-arm carry.} One object held by both grippers and moved without
dropping or tilting it. Currently not performable, for a reason given in the
next chapter: the object presents more width to the gripper than the gripper
opens.

\emph{T4, handing to the wearer.} The robot picks something up and offers it
to the person carrying it. It is the only task in which the machine serves the
wearer rather than working past them, and it is the one where ownership and
trust have something to attach to.

\section{What is measured}

\begin{table}[htbp]
  \centering
  \small
  \caption[Measures]{Everything recorded, by whom it is answered and when.}
  \label{tab:measures}
  \begin{tabular}{@{}p{0.30\linewidth}p{0.30\linewidth}p{0.32\linewidth}@{}}
    \toprule
    \textbf{Measure} & \textbf{Source} & \textbf{When} \\
    \midrule
    Completion time & Logged automatically & Every trial \\
    Path efficiency & Logged automatically & Every trial \\
    Success or failure & Logged automatically & Every trial \\
    Number of clutch re-engagements & Logged automatically & Every trial \\
    Closest approach to the wearer & Logged automatically & Every trial \\
    Emergency stops and their cause & Logged automatically & Every trial \\
    Task load index~\cite{hart1988nasatlx} & Both participants & After each
      block \\
    Physical exertion~\cite{borg1982perceived} & Wearer & After each block \\
    Trust in the system & Both participants & After each block \\
    Sense of ownership~\cite{longo2008embodiment} & Wearer & End of session \\
    Free comment & Both participants & Debrief \\
    \bottomrule
  \end{tabular}
\end{table}

Exertion is asked of the wearer between every block rather than once at the
end. Fatigue in a task like this rises steadily, and asking once at the end
turns it from something measured into something guessed at.

\section{How a session runs}

\begin{figure}[htbp]
  \centering
  % What happens in one two-hour session.
\begin{tikzpicture}[x=1mm, y=1mm, font=\scriptsize]

  \foreach \i/\t/\lbl/\who in {
      0/{0:00}/{Consent, separately}/{both, in different rooms},
      1/{0:20}/{Fit the harness, check the fit}/{wearer},
      2/{0:35}/{Familiarisation, no task}/{driver},
      3/{0:50}/{Prove the emergency stop}/{all three present},
      4/{1:00}/{Block 1: one condition, all tasks}/{both},
      5/{1:20}/{Rest, questionnaires, exertion rating}/{both},
      6/{1:30}/{Block 2}/{both},
      7/{1:50}/{Rest and questionnaires}/{both},
      8/{2:00}/{Block 3, debrief}/{both}}
  {
    \pgfmathsetmacro{\y}{-\i*9}
    \node[anchor=east, font=\scriptsize\ttfamily] at (0,\y) {\t};
    \ifnum\i=3
      \node[guard, anchor=west, text width=54mm, minimum height=6mm] at (4,\y) {\lbl};
    \else
      \ifnum\i=5 \node[person, anchor=west, text width=54mm, minimum height=6mm] at (4,\y) {\lbl};
      \else \ifnum\i=7 \node[person, anchor=west, text width=54mm, minimum height=6mm] at (4,\y) {\lbl};
      \else \node[stage, anchor=west, text width=54mm, minimum height=6mm] at (4,\y) {\lbl};
      \fi \fi
    \fi
    \node[anchor=west, font=\tiny\itshape, text width=34mm] at (62,\y) {\who};
  }

  \draw[decorate, decoration={brace, amplitude=4pt}] (100,-36) -- (100,-72)
        node[midway, right=5pt, font=\tiny\itshape, text width=26mm]
        {The wearer carries the load for no more than twenty minutes at a
         time, with a rest between blocks};

\end{tikzpicture}

  \caption[A session]{One session, about two hours. The red row is not
  optional: the emergency stop is demonstrated and tested with both
  participants watching before any task begins, and the wearer is shown how to
  stop the arms themselves.}
  \label{fig:session}
\end{figure}

Three people are present. The two participants, and a third person whose only
job is to hold the stop and watch the wearer rather than the screen. That
third role exists because the driver is concentrating on the task and the
wearer cannot see behind their own shoulders.

Consent is taken separately and in different rooms. The two participants may
know each other, and a wearer who feels watched by the person who brought them
is not consenting freely.

\section{Analysis}

For each measure, the comparison is across conditions within a pair. Where the
assumptions hold, a repeated-measures analysis of variance; where they do not,
its rank-based equivalent. Post-hoc comparisons are corrected for multiple
tests, and both the raw and the corrected values are reported side by side so
that a reader can see what the correction did.

Effect sizes are reported with confidence intervals, and the intervals rather
than the significance are what the discussion will be written around. With
sixteen pairs and three trials each, the study is powered to detect a large
effect on completion time and is not powered to detect a small one. That is
stated in advance rather than discovered afterwards.

Trials that were invalidated by a fault are excluded, and the rule for
exclusion is written into the analysis code rather than applied by hand. If
more than half a session's trials are invalid, the analysis stops and says so,
because a session that fails silently and produces empty averages is
indistinguishable from a null result.

\section{Ethics and safety}

The application is written and the study has not been submitted for approval,
because a study cannot honestly be submitted while its baseline condition
depends on sensor channels that do not work.

The largest single risk is easy to state: a seventeen-kilogram machine mounted
behind a seated person, moving under partly automatic control, within reach of
their head. Everything else follows from that.

\begin{table}[htbp]
  \centering
  \small
  \caption[Risks and what is done about them]{The main risks to participants
  and the measures in place. Software measures are the ones described in
  \cref{ch:safety}; the rest are procedural.}
  \label{tab:risks}
  \begin{tabular}{@{}p{0.34\linewidth}p{0.58\linewidth}@{}}
    \toprule
    \textbf{Risk} & \textbf{What is done} \\
    \midrule
    Contact with the wearer & A \SI{150}{\milli\metre} margin enforced along
      the whole path, measured geometrically, with the arm holding rather than
      slowing when it is reached \\
    The wearer cannot stop the arms & A stop within the wearer's own reach, and
      a third person holding another \\
    Load on the wearer's back & Twenty minutes carrying at a time, seated where
      possible, exertion rated between blocks, and the session ended on request
      without explanation \\
    Fatigue confounding the results & Condition order balanced; exertion
      recorded as a covariate rather than assumed constant \\
    A participant feeling unable to withdraw & Consent taken separately, and
      the right to stop restated at every break \\
    Identifying data & Refused in code: the routine that writes the session
      record raises an error if given a name, an email address, a date of
      birth, an address or a telephone number \\
    \bottomrule
  \end{tabular}
\end{table}

Two things are worth saying about the last row. It is enforced in the program
rather than in a policy document, because a policy document does not stop a
keystroke. And video recording, which would be useful, is not part of the
protocol: a video of somebody wearing this rig is not anonymous, and the
analysis does not need one.
